\section{Worked application templates}
\label{sec:applications}

This section illustrates how the framework changes proof shape.  The examples
are templates rather than claims that every listed deep theorem has now been
formalized.

\subsection{A value query with a dead branch witness}

Suppose a finite algebra procedure produces a certificate containing a root
cluster \(C\), an ordered branch label \(\ell\), a local trace \(t\), and an
error bound \(e\).  The public query asks only for a symmetric observable
\(s(C)\) to precision \(p\).  The resource rules may be:
\begin{align*}
  \text{query:value}(p)&\Rightarrow \Prec(s)\ge p,\\
  \Prec(s)\ge p&\Rightarrow \Prec(C)\ge p+c,\\
  \text{query:labeled-root}&\Rightarrow \Cert(\ell)=1,
       \Guard(\text{separation})=1,\\
  \text{query:replay}&\Rightarrow \Cert(t)=1.
\end{align*}
For the value-only seed, the least closure leaves \(\ell\) and \(t\) dead and
never requests separation.  The certificate producer is opened locally; the
symmetric-observer theorem derives the value; existential elimination seals
the full certificate.  The endpoint is stronger and more stable than one that
exports an arbitrary ordered root.

If the query changes to a labeled root, query monotonicity activates the label,
separation guard, and named export.  The stronger theorem visibly has stronger
hypotheses.

\subsection{A cyclic precision closure}

The generated demonstration uses coordinates
\[
 \Prec(\text{root})\in\{0,\ldots,6\},
 \quad
 \Prec(\text{integral})\in\{0,\ldots,6\},
\]
with rules
\begin{align*}
  \text{value query}&\Rightarrow \Prec(\text{root})\ge3,\\
  \Prec(\text{root})\ge1&\Rightarrow
       \Prec(\text{integral})\ge2,
       \Guard(\text{denominator})=1,\\
  \Prec(\text{integral})\ge2&\Rightarrow
       \Prec(\text{root})\ge5,\\
  \Prec(\text{root})\ge5&\Rightarrow
       \Prec(\text{integral})\ge6,
       \Res(\text{tail modulus})=1.
\end{align*}
The least closure is
\[
  \Prec(\text{root})=5,
  \quad \Prec(\text{integral})=6,
  \quad \Guard(\text{denominator})=1,
  \quad \Res(\text{tail modulus})=1.
\]
Branch and trace certificate coordinates remain zero; so do a separation guard
and an unrelated height-bound residual.  The cycle is solved entirely in the
finite metalevel lattice.  The object proof is a straight schedule at the
chosen precision levels.

\subsection{Bounded alternation}

Consider the finite statement
\[
  \forall i<I\;\exists j<J\;R(i,j).
\]
A direct nested proof can be awkward because the existential witness depends on
\(i\).  DARC generates a strategy certificate coding the function
\(i\mapsto j_i\), together with a checker asserting
\(j_i<J\) and \(R(i,j_i)\) for all \(i<I\).  The checker is primitive
recursive and its validity reflects the original bounded formula.  Three
public endpoints are possible:
\begin{enumerate}
  \item seal the strategy and prove only the bounded formula;
  \item observe a stable aggregate, such as a finite sum of costs;
  \item name the least valid strategy code under a finite code bound.
\end{enumerate}
The reference artifact instantiates this with parity and verifies the finite
certificate.

\subsection{A deterministic recurrence}

For the loop
\[
  (n,a)\mapsto(n+1,a+2n+1),
\]
the invariant \(a=n^2\) is polynomial and quantifier-free.  DARC generates the
trace functions and applies the semiring proof of
\((n+1)^2=n^2+2n+1\).  No trace list is public.  A replay query would activate a
trace code, but the theorem of the recurrence itself does not need one.

This pattern covers many existing checked derivations: finite sums, powers,
recurrences, the Euclidean algorithm, and simultaneous recursions.  DARC's
contribution is to package them as loop nodes with explicit transfer summaries
and certificate scope.

\subsection{Finite compatible approximation clouds}

Suppose each precision level \(h\le H\) has a finite set of approximation
certificates \(X_h\), and \(E_h(x,y)\) means that \(y\) refines \(x\).  Local
nonemptiness of every \(X_h\) is insufficient: independently least elements
may be incompatible.  DARC computes viability backward from terminal
accuracy, prunes dead certificates, and only then chooses least successors.

If compatible edges preserve an observable within error \(\varepsilon_h\), the
public endpoint is the root observation with total finite error
\(\sum_{h<H}\varepsilon_h\).  The path is sealed.  To pass from all finite
horizons to a regular name, the plan must activate an extension function and a
tail-modulus producer.

\subsection{Special functions and analytic packets}

A special-function packet typically mixes four levels:
\begin{enumerate}
  \item exact identities of finite products, recurrences, or coefficients;
  \item formal power-series identities at a fixed truncation;
  \item effective convergence inside a guarded domain;
  \item continuation, branch selection, asymptotics, or special values.
\end{enumerate}
DARC separates these levels.  Exact finite identities often compile with no
residual.  A truncation theorem requests a finite coefficient bound.  An
analytic value query activates a convergence or tail modulus and domain
margins.  A branch-specific query additionally activates a path and separation
guards.  A continuation theorem beyond the formalized domain remains residual.

This prevents an exact coefficient proof from being overstated as a complete
theorem about a global analytic function, while maximizing the portion that
can be checked in \(\iPRA\).

\subsection{Diophantine enumeration}

Suppose \(R(x,y)=0\) is a concrete polynomial or primitive-recursive relation.
The verifier is straightforward.  Given a height \(H\), bounded search can
enumerate all solutions with code below \(H\), prove soundness of the list, and
prove completeness below \(H\).  To claim global completeness, the plan needs
\[
  R(x,y)=0\Longrightarrow \operatorname{height}(x,y)<H.
\]
If no effective height producer is known, DARC emits a conditional theorem
with that formula as the sole deep residual.  This is more informative than
marking the whole application ``unformalizable'': the verifier, enumeration,
and downstream finite reasoning can still be completed.

\subsection{Geometry and finite nets}

For a compact-domain geometric inequality, a DARC plan can contain:
\begin{enumerate}
  \item a rational mesh generator at scale \(p\);
  \item a finite-net coverage producer;
  \item a Lipschitz or continuity modulus;
  \item exact checking of the inequality with slack on every mesh point;
  \item transfer from mesh slack to the full domain.
\end{enumerate}
The requested final slack determines mesh precision.  Degenerate configurations
activate separate guards or case nodes.  If the coverage or modulus is not yet
formalized, it remains an exact residual rather than being supplied by
abstract compactness.

\subsection{Busy Beaver and exact coded endpoints}

Machine and consistency applications demand exact provenance.  A plan should
bind the literal table, blank-start convention, proof predicate, and object
theory at producer nodes.  Generic arithmetic lemmas can be imported through
ordinary open sequents; a final bridge node translates to the published
Foundation endpoint.  DARC does not identify two machines because their
external simulations agree, nor does it replace a conditional object theorem
by a stronger semantic statement.  The endpoint inventory remains the source
of truth.
